\section{Серверная сторона}
\label{sec:server}

\textbf{Измеренный сервер.} Базовая кривая получена на RTX PRO 6000 Blackwell Workstation
с 96~ГБ и лимитом мощности 600~Вт: llama.cpp \texttt{c21284cd}, CUDA 12.8, Qwen3.8-27B Q4\_K\_M,
кэш \texttt{q8\_0}, Flash Attention и полная выгрузка слоёв на GPU, контрольные точки отключены.
Повторный запуск воспроизвёл время на токен с максимальным расхождением 0.355\%.
На глубине 16 тысяч порции в 512, 2000 и 10\,000 токенов дали 3509, 3547 и 3435 токенов/с;
на глубине 249 тысяч, соответственно 898, 905 и 971 токен/с.
Генерация на концах диапазона составила 61.8 и 34.6 токена/с.

С настройками контрольных точек по умолчанию те же порции дали 1600, 2876 и 3263 токена/с
на глубине 16 тысяч, 742, 946 и 973 токена/с на глубине 249 тысяч; генерация не изменилась.
Издержки контрольных точек заметнее на малых порциях. Более ранняя серия на той же модели
карты и версии движка дала для порции 2000 токенов 2147 и 533 токена/с на концах диапазона,
генерация была на 15--30\% медленнее. Проверка контрольных точек причину раннего замедления
не объяснила. Расчёты сохраняют раннюю серию и режим контрольных точек как отдельные случаи чувствительности.

Эффективная доля плотного пика FP8 равна отношению расчётных операций модели к измеренному
времени prefill и паспортному пику GPU. В базовой кривой эта доля составляет 14.1--21.1\%
по глубинам против 8--13\% в ранней серии. Для Qwen3.8-27B доля вычисляется как
$(54\cdot10^9 + 393216L)/[(a+bL)P_{\mathrm{FP8}}]$, где $a+bL$ аппроксимирует секунды на токен.
Эта величина характеризует эффективную скорость модельных операций; загрузка GPU по телеметрии
и фактический формат арифметики Q4-ядер из этой величины не следуют.
Перенос кривой на GLM остаётся оценкой.

Отдельный тест \texttt{llama-bench} показал ускорение новой версии \texttt{dde879bc} на 9--20\%
относительно \texttt{c21284cd}; подбор размера микропакета добавил 2--5\%.
Для порции 4096 токенов с нулевой глубины исходная версия с микропакетом 4096 дала 3950 токенов/с,
новая с микропакетом 512 дала 4793 токена/с. Базовые экономические расчёты сохраняют
воспроизведённую кривую исходной версии.

\textbf{Серверный FP8.} vLLM 0.32.0 с Qwen3.8-27B-FP8 на RTX PRO 6000 показал
аппроксимацию холодного prefill: она даёт для порции 10\,000 токенов 8122 токена/с на глубине 16 тысяч
и 2209 на глубине 249 тысяч. Скорости относятся к конечной порции по аппроксимации времени на токен.
Кэширование префикса работает в режиме Mamba \texttt{align} с блоками по 784 токена.
Для добавочных порций в 2000 и 10\,000 токенов полезные скорости составили 5370/7144 токена/с
на глубине 16 тысяч, 3162/4598 на 65 тысячах, 2177/3130 на 131 тысяче и 1258/1962 на 249 тысячах.
Скорости считают новые токены порции; непрокэшированный остаток блока также требует вычислений.
Режимы NVFP4 и FP8 для кэша внимания на PRO 6000 не запустились.
Эти скорости FP8 относятся только к серверному исполнению.

\textbf{Измеренная B300.} B300 SXM6 AC с 275~ГБ и лимитом мощности 1100~Вт измерена на CUDA 12.9
комплектом 1.6.5: llama.cpp \texttt{dde879bc}, сборка \texttt{103a-real}, Qwen3.8-27B Q4\_K\_M,
кэш \texttt{q8\_0}, контрольные точки отключены. Таблица~\ref{tab:b300-llama} показывает лестницу глубин.
Для порции 2000 токенов скорость равна 1659 токенам/с на глубине 16 тысяч и 892 на 249 тысячах,
против 3547 и 905 у PRO 6000. На этих глубинах prefill B300 примерно равен или уступает PRO 6000.
Сравнение использует разные версии llama.cpp: \texttt{dde879bc} и \texttt{c21284cd} соответственно.
Масштабирование скорости по отношению паспортных пиков не описывает этот замер.

\begin{table}[t]
\centering
\caption{B300, llama.cpp \texttt{dde879bc}, комплект 1.6.5. Все скорости в токенах/с,
глубина в тысячах токенов; размер порции задан числом новых токенов.}
\label{tab:b300-llama}
\footnotesize
\begin{tabular}{rrrrrr}
\toprule
Глубина & Сегмент prefill & Порция 512 & Порция 2000 & Порция 10000 & Генерация \\
\midrule
16 & 1757 & 1658 & 1659 & 1663 & 79.6 \\
50 & 1580 & 1456 & 1475 & 1485 & 67.6 \\
100 & 1388 & 1232 & 1267 & 1284 & 54.8 \\
150 & 1212 & 1066 & 1110 & 1131 & 46.3 \\
200 & 1074 & 941 & 990 & 1010 & 40.1 \\
249 & 966 & 844 & 892 & 915 & 35.4 \\
\bottomrule
\end{tabular}
\end{table}

\textbf{FP8 на B300.} vLLM 0.32.0 с Qwen3.8-27B-FP8 даёт скорости очередного токена по аппроксимации
холодного prefill, приведённые в таблице~\ref{tab:b300-vllm}. Они характеризуют локальную скорость
при указанной глубине; скорость конечной порции получается интегрированием времени на токен.
Для порции 10\,000 токенов на глубине 16 тысяч аппроксимация даёт 23\,151 токен/с
против 1663 у llama.cpp на той же B300. Кэш FP8 на B300 запустился.
Кэшированные порции в 2000/10\,000 токенов с обычным кэшем и блоками по 784 токена дали
11\,638/20\,140 на глубине 16 тысяч, 8755/14\,768 на 65 тысячах, 7251/11\,185 на 131 тысяче
и 4134/6987 токенов/с на 249 тысячах. С кэшем FP8 и блоками по 1568 токенов на глубине
16 тысяч получено 13\,761/20\,961 токен/с. Все эти скорости относятся только к серверному исполнению.

\begin{table}[t]
\centering
\caption{B300, vLLM 0.32.0 FP8: скорости очередного токена по аппроксимации холодного prefill,
токенов/с, только сервер. Глубина в тысячах токенов.}
\label{tab:b300-vllm}
\footnotesize
\begin{tabular}{rrr}
\toprule
Глубина & Обычный кэш & Кэш FP8 \\
\midrule
16 & 24031 & 24841 \\
50 & 19096 & 21846 \\
100 & 14667 & 18555 \\
249 & 8673 & 12806 \\
\bottomrule
\end{tabular}
\end{table}

\textbf{Первый запуск моста состояния.} B300 с vLLM 0.32.0 FP8 экспортировала состояние на
784 токенах; клиент llama.cpp на той же карте сравнивался со своим локальным prefill на 64 позициях.
Клиент BF16 получил совпадение top-1 в 64/64 позициях и средний KL 0.00055 нат,
клиент Q4\_K\_M совпал в 64/64 позициях при среднем KL 0.0027 нат.
Свободная генерация не разошлась на 64 токенах в обоих случаях. Это гетерогенная некалиброванная
проверка: серверные веса FP8 отличаются от клиентских GGUF. Проверка моста измеряет приближённую эквивалентность
на этих позициях и не устанавливает побитовое совпадение или скорость полного гибридного вызова.
На глубине 31\,360 токенов средний KL равен 0.0031 нат у BF16 и 0.0061 нат у Q4\_K\_M, top-1 совпал в 64/64 позициях у обоих клиентов.
Свободная генерация не разошлась за 64 токена. Проверка остаётся гетерогенной и некалиброванной;
код выхода проверяющей программы 5 сохраняет этот статус.
\textbf{Сквозная проверка узла vLLM.} Адаптер на B300 связал vLLM 0.32.0 FP8,
стандартный прокси и клиент llama.cpp Q4\_K\_M на той же карте; все вставки завершились.
На глубине 8 тысяч гибрид дал KL 0.0055 нат, top-1 63/64, отношение к собственному шуму разбивки клиента 0.87.
На глубине 32 тысячи получены KL 0.0256 нат, top-1 60/64 и отношение 2.38;
свободная генерация разошлась на девятом токене. Среднее время до первого токена по восьми вызовам
равно 2.72 секунды локально и 4.15 секунды при оффлоаде. Одна GPU делает этот запуск функциональной проверкой.
Проверка гетерогенная и некалиброванная, FP8 на узле и Q4 на клиенте; код выхода 5 сохраняет этот статус.

Точная общая версия прототипа \texttt{9348ecd6} после PR \#16 повторно проверена на RTX PRO 6000:
узел vLLM 0.32.0 FP8 и клиент Q4\_K\_M работают на одной карте. Сбор состояния и побайтовая
проверка CUDA прошли. Среднее время до первого токена по восьми вызовам составило 1.571 секунды
локально и 4.497 секунды при оффлоаде, генерация 73.6--74.5 токена/с; проверка функциональная.
При 8 тысячах получены KL 0.0202 нат, top-1 61/64 и первое расхождение на восьмом токене;
собственный шум разбивки клиента 0.0053--0.0099 нат, top-1 62--63/64, отношение 2.03.
На B300 при 8 тысячах измерен KL 0.0055 нат; причина различия между картами не установлена.
При 32 тысячах PRO 6000 даёт KL 0.0233 нат, top-1 60/64 и расхождение на девятом токене;
шум разбивки 0.0022--0.0030 нат при top-1 63--64/64, отношение 7.83.
Дельта пришла пятью частями; добавочная цепочка с концами 7840, 31\,360 и 32\,144 дала те же числа.
Повтор одинаковой разбивки побитово точен, KL равен нулю. Код выхода 5 сохраняет статус
гетерогенной некалиброванной проверки. Измеренные времена B300 для точной проверенной версии
\texttt{c2f47472} равны 2.69 секунды локально и 4.04 секунды при оффлоаде.

Экономические оценки B300 используют измеренную кривую llama.cpp; цены vLLM обозначены как серверные без квалификации полного гибридного вызова.
Модель учитывает вычисление новых токенов и ответов клиента, подъём префикса и запись новых позиций.
В сравнении собственных серверов активная генерация занимает память GPU, а сохранённые сессии
оплачены отдельным пулом хранения. Оба маршрута могут вытеснять неактивные сессии
и выбирают HBM, оперативную память или NVMe при одинаковом бюджете задержки.

\textbf{Хранение сессий.} Сессия закреплена за узлом, который держит её кэш в памяти GPU: по нашей оценке на H100 с 80~ГБ после весов
помещается около 16 сессий по 100 тысяч токенов. Неактивные сессии узел вытесняет в оперативную память
сервера, откуда 4~ГБ возвращаются за 0.1--0.2 секунды, затем на локальный NVMe (0.3--1.3 секунды) и в объектное
хранилище, которое служит только для восстановления. Даже самое медленное восстановление дешевле повторного
чтения всего контекста: на H100 это по нашей оценке от 13 до 33 секунд на глубине 172 тысячи токенов, близкой к средней глубине профиля (164 тысячи). Поэтому серверу нужно железо под своим
управлением, своё или арендованное целиком: посекундные бессерверные GPU не хранят кэш сессий между вызовами.

Редкие обращения не освобождают копию в модели: расчёт оплачивает четыре сохраняемые сессии на пользователя
полный календарный месяц. Политика протокола держит копии в оперативной памяти в рабочие часы и на NVMe
в остальное время; осторожная политика держит копии в оперативной памяти весь месяц.
Это заданные сценарии хранения. Оптимальный срок удаления по неактивности не рассчитан и не измерен.

\textbf{Сколько людей держит одна карта.} При рабочем дне в 8 часов и загрузке 70\% одна RTX PRO 6000 держит prefill
для 33 тяжёлых пользователей при доле открытой модели 20\%, для 13 при доле 50\% и для 6, если на открытую модель
уходит вся их работа. Одна B300 с llama.cpp держит 27, 11 и 5 человек. Эти числа учитывают подъём и запись копий сессий по
политике протокола; одни вычисления дали бы 35, 14 и 7 на RTX PRO 6000 и 28, 11 и 5 на B300, а осторожная
политика хранения даёт 31, 12, 6 и 26, 10, 5. При доле 20\% одной карты среднего класса хватает команде в несколько
десятков человек, при доле 50\% примерно тринадцати.
